\section{Demo final}
\label{sec:demo-final}

Como cierre del proyecto se preparan dos demostraciones equivalentes del pipeline completo. La primera utiliza la ruta local basada en XTTS, mientras que la segunda utiliza la integración con ElevenLabs para la síntesis de voz. Ambas demos mantienen la misma estructura de procesamiento: detección, tracking, identificación de jugadores, clasificación de acciones, generación de comentarios y ensamblado audiovisual final.

La demo con XTTS representa la configuración local por defecto del sistema. Su principal interés es mostrar una ejecución sin dependencia de servicios externos y con menor coste por uso. El video de prueba puede consultarse en el siguiente enlace:

\begin{center}
\fbox{\href{https://github.com/hectorgarciaa/narrador-futbol/raw/final/demo_final/demo_final_xtts.mp4}{\texttt{demo\_final\_xtts.mp4}}}
\end{center}

La segunda demo utiliza ElevenLabs como backend de voz. Esta configuración se reserva para una demostración de mayor calidad perceptiva, ya que las voces diseñadas mediante Voice Design resultan más naturales y expresivas, aunque introducen dependencia de API, consumo de créditos y una latencia inicial asociada al servicio externo. El video de prueba puede consultarse en:

\begin{center}
\fbox{\href{https://github.com/hectorgarciaa/narrador-futbol/raw/final/demo_final/demo_final_elevenlabs.mp4}{\texttt{demo\_final\_elevenlabs.mp4}}}
\end{center}

Además de los videos con narración, la interfaz genera una salida visual específica de PathCRF. Este video no necesita contener audio: su función es mostrar la clasificación de acciones sobre el partido, con las aristas y eventos que alimentan posteriormente al sistema de comentarios. El video de prueba correspondiente puede consultarse en:

\begin{center}
\fbox{\href{https://github.com/hectorgarciaa/narrador-futbol/raw/final/demo_final/demo_final_pathcrf.mp4}{\texttt{demo\_final\_pathcrf.mp4}}}
\end{center}

Los videos enlazados en esta sección se han dejado con nombres estables para poder sustituirlos por las versiones finales sin modificar la memoria.
